\chapter{Counting and Discrete Reasoning}\label{ch:counting}

\begin{center}\itshape
How many ways can every guest receive the wrong hat?\\
The difficulty is to count overlapping restrictions exactly once.
\end{center}

Chapter~\ref{ch:sums} counted by adding, and it ended with a count that is not a sum of disjoint pieces. The binomial coefficient counts the subsets of an $n$-element set. A hat check asks for something thinner: the permutations in which nobody receives their own hat. If you list the arrangements that disappoint guest~$1$, and then the arrangements that disappoint guest~$2$, some arrangements are written twice. This chapter is about counting each object once. Inclusion--exclusion is the correction for an overlap. Lattice paths are the same binomial numbers, read as walks. A generating function then stores a whole list of counts as one polynomial.

Three guests are enough to see the problem. Write a permutation as the word whose $i$th letter is the hat received by guest~$i$. There are $3!=6$ words.
\[
123,\quad 132,\quad 213,\quad 321,\quad 231,\quad 312.
\]
In $123$ every guest is fixed. In $132$ guest~$1$ is fixed, in $213$ guest~$3$ is fixed, and in $321$ guest~$2$ is fixed. The words $231$ and $312$ fix nobody. So exactly two of the six arrangements are derangements. Figure~\ref{fig:hats} is that list. Six arrangements are not a formula. The rest of the chapter builds the rules that still work for $n$ guests.

\begin{figure}[ht]
\centering
\fig{ch02-hats}
\caption{The six permutations of three hats, written as the hats received by guests $1$, $2$ and $3$. A rust digit is a guest who received their own hat. Four of the six words contain a rust digit. The two charcoal words $231$ and $312$ contain none, so $D_3=2$.}
\label{fig:hats}
\end{figure}

The same question is old. In his \emph{Essay d'analyse sur les jeux de hazard} (1708), Pierre R\'emond de Montmort asked it for a pack of thirteen cards: how many shuffles leave no card in its original place? He called this the \emph{probl\`eme des rencontres}. The second edition, of $1713$, prints his correspondence with Nicolaus Bernoulli on that problem. Euler later treated the same count. The sources at the end of the chapter give both references. We count the arrangements in this chapter. What the count says about chance waits for Chapter~\ref{ch:probability}.

\section{Counting, from scratch}

A counting argument has to name its objects, and it has to say why each object is counted once. Most of the decisions happen before any factorial is written. Sometimes the complement is easier than the set itself. Sometimes the objects split into disjoint cases. Sometimes the same object has a second description, and the second description is the one we know how to count.

\begin{reminder}[the three rules of counting]
For finite sets, $|A\cup B|=|A|+|B|$ when $A\cap B=\emptyset$. This is the \emph{sum rule}.
$|A\times B|=|A|\,|B|$. This is the \emph{product rule}. It extends to a choice made in $k$ stages, with $n_1,\dots,n_k$ options at the stages, even when the list of options depends on earlier choices, provided the \emph{number} of options does not. Finally $|A|=|B|$ whenever there is a bijection $A\to B$. This is the \emph{bijection rule}.
Almost every count in this chapter is one of these three. The usual mistake is to count an object twice. The cure is to check that the stages of a product-rule count determine the object, and are determined by it.
\end{reminder}

Each rule is a small bijection in disguise. For the sum rule, list the elements of the disjoint sets one list after the other. No element is repeated, and none is omitted, because an element of the union lies in exactly one of the sets. For the product rule, group the pairs in $A\times B$ by their first coordinate. There are $|A|$ groups, they are disjoint, and each group has $|B|$ pairs, so the product rule is the sum rule applied $|A|$ times. Repeating the same grouping gives the stage-by-stage form. A bijection carries a list of one set onto a list of the other, so the lengths agree.

Here is the complement, on a set small enough to say what was avoided. How many five-letter words on a $26$-letter alphabet contain at least one vowel? A direct count splits into words with one vowel, two vowels, and so on. The complement does not. There are $26^5$ words in all, and $21^5$ words that use only the $21$ consonants. The difference is the count we want. The powers are short to recompute: $26^2=676$ and $676^2=456{,}976$, so $26^5=456{,}976\cdot 26=11{,}881{,}376$. Likewise $21^2=441$ and $441^2=194{,}481$, so $21^5=194{,}481\cdot 21=4{,}084{,}101$. Therefore
\[
26^5-21^5=11{,}881{,}376-4{,}084{,}101=7{,}797{,}275.
\]

The next two counts are bijections, and each one has a check you can draw. How many ways are there to seat $n$ people at a round table, if two seatings are the same whenever one is a rotation of the other? For $n=3$, label the people $A,B,C$. The clockwise orders are $ABC$ and $ACB$. A rotation of $ABC$ is still $ABC$, so there are $2$ seatings. The same count is $3!/3=2$, or $(3-1)!=2$. For $n=8$ the same reasoning gives $8!/8=7!=5040$. One proof of the division: the $8!$ linear seatings fall into groups of $8$ rotations, and the groups are the round seatings. Another proof: fix one person's place, which kills the rotational symmetry, and arrange the other $7$ people in the remaining seats in $7!$ ways. The two answers agree. Reflections have not been identified. A seating and its mirror image are different under this count.

How many diagonals has a convex $n$-gon? A diagonal joins two vertices and is not a side. There are $\binom n2$ pairs of vertices and $n$ sides, so the count is $\binom n2-n$. For a square, $n=4$ and $\binom42-4=2$, the two diagonals you can draw. In general
\[
\binom n2-n=\frac{n(n-1)}{2}-n=\frac{n(n-3)}{2}.
\]
A pentagon has $5\cdot 2/2=5$ diagonals.

One more case split, because it returns in the workshop. Let $f_n$ be the number of subsets of $\{1,\dots,n\}$ that contain no two consecutive integers. The empty subset is allowed. It is one subset, and the condition does not reject it. For $n=1$ the subsets are $\emptyset$ and $\{1\}$, so $f_1=2$. For $n=2$ they are $\emptyset$, $\{1\}$ and $\{2\}$, so $f_2=3$. The pair $\{1,2\}$ is forbidden.

Now take $n\ge 3$. A good subset either avoids $n$, or it contains $n$. If it avoids $n$, it is a good subset of $\{1,\dots,n-1\}$, and there are $f_{n-1}$ of those. If it contains $n$, it cannot contain $n-1$, so the rest is a good subset of $\{1,\dots,n-2\}$, and there are $f_{n-2}$ of those. The two cases are disjoint, and every good subset falls into one of them. Thus $f_n=f_{n-1}+f_{n-2}$. The next value is $f_3=3+2=5$. The five subsets of $\{1,2,3\}$ are $\emptyset$, $\{1\}$, $\{2\}$, $\{3\}$ and $\{1,3\}$. The Fibonacci numbers of Exercise~\ref{ex:fibgcd} obey the same recurrence, with $F_0=0$, $F_1=1$, $F_3=2$ and $F_4=3$. The workshop matches the two sequences and finds a closed form.

\begin{proposition}[basic counts]\label{prop:counts}
Let $|E|=n\ge1$ and $|F|=k\ge0$.
\begin{enumerate}[label=(\alph*)]
\item There are $n^k$ maps $F\to E$, $\frac{n!}{(n-k)!}$ injections for $k\le n$ (and none for $k>n$), and $n!$ permutations of $E$.
\item There are $\binom nk$ subsets of size $k$ and $\binom{n+k-1}k$ multisets of size $k$ (``stars and bars'').
\item \textbf{Pigeonhole}: if $m\ge0$ objects go into $n$ boxes, some box holds at least $\lceil m/n\rceil$ of them, where $\lceil x\rceil$ is the least integer greater than or equal to $x$.
\end{enumerate}
\end{proposition}
\begin{proof}
\emph{Maps, injections and permutations.} Enumerate $F=\{f_1,\dots,f_k\}$. A map chooses an image in $E$ for each $f_i$, and the choices are independent, so there are $n$ options at each of the $k$ stages. The product rule gives $n^k$. For $k=0$ the product of no factors is $1$: there is one map from the empty set to $E$, the empty function. An injection must choose distinct images. There are $n$ choices for the image of $f_1$, then $n-1$ remaining choices for $f_2$, and so on, down to $n-k+1$ choices for $f_k$, provided $k\le n$. The product is
\[
n(n-1)\cdots(n-k+1)=\frac{n!}{(n-k)!}.
\]
Chapter~\ref{ch:sums} records $0!=1$, so the case $k=0$ is $n!/n!=1$, the one empty injection. If $k>n$, the choices run out: after $n$ distinct images have been chosen, the next element has nowhere left to go, so there are no injections. A permutation of $E$ is a bijection $E\to E$, which is an injection from a set of size $n$ to itself, so there are $n!$ of them.

\emph{Subsets and multisets.} The subset count is Proposition~\ref{prop:binom} in Chapter~\ref{ch:sums}. For multisets, let $E=\{e_1,\dots,e_n\}$. A multiset of size $k$ is a choice of multiplicities $m_1+\dots+m_n=k$ with each $m_i\in\N$. Encode it as a word of $k$ stars and $n-1$ bars: $m_1$ stars, a bar, $m_2$ stars, a bar, and so on, ending with $m_n$ stars. A bar may stand next to another bar, and that encodes a zero multiplicity. This is a bijection between the multisets and the words of length $n+k-1$ with exactly $k$ stars. Such a word is determined by the set of positions of its stars, a $k$-element subset of $\{1,\dots,n+k-1\}$. There are $\binom{n+k-1}k$ of them.

\emph{Pigeonhole.} Let $c=\lceil m/n\rceil$. Then $c-1<m/n$. Indeed $c$ is an integer, $c\ge m/n$, and $c-1$ is a strictly smaller integer, so $c-1$ cannot be $\ge m/n$: otherwise $c$ would not be the least integer with that property. If every box held at most $c-1$ objects, the total number of objects would be at most $n(c-1)$. Since $c-1<m/n$ and $n>0$, that total is strictly less than $m$. So the supposition is false, and some box holds at least $c$ objects.
\end{proof}

Check the three parts on sets you can list. Take $n=2$ and $k=2$. There are $2^2=4$ maps from a $2$-element set to a $2$-element set, $2!/(2-2)! = 2$ injections, and $2!=2$ permutations. The $2$-element subsets of a $4$-element set number $\binom42=6$. The multisets of size $2$ on $\{a,b\}$ are $\{a,a\}$, $\{a,b\}$ and $\{b,b\}$, three of them, and $\binom{2+2-1}{2}=3$. The three words are two stars then a bar, a star then a bar then a star, and a bar then two stars. For the pigeonhole principle, put $5$ books on $3$ shelves. Then $\lceil 5/3\rceil=2$, so some shelf holds at least two books. If every shelf held at most one, there would be at most $3$ books.

The pigeonhole principle is the remark that a finite set cannot be spread thinner than its size allows. Its strength is that it proves existence without exhibiting the object. In $1842$ Dirichlet used it to prove that if $\alpha$ is irrational, then infinitely many rationals $p/q$, with $q>0$, satisfy $\bigl|\alpha-p/q\bigr|<1/q^2$. The argument places the fractional parts of $0,\alpha,2\alpha,\dots,N\alpha$ into $N$ equal intervals. A fractional part is a real number minus its floor, and the floor of a general real number is defined in Chapter~\ref{ch:sequences}. This chapter does not use that floor, so the approximation argument is not given here. Exercises~\ref{ex:pigeon} and~\ref{ex:erdos} show what a careful choice of boxes can do with nothing but the principle above.

The counts in the proposition are exact when each object meets the condition in one way. The hat check is the opposite situation. An arrangement can disappoint several guests at once, and a plain sum counts that arrangement several times.

\section{Inclusion--exclusion}

Begin with a union of two sets, and write every element down. Among the integers $1,\dots,10$, let $A$ be the multiples of $2$ and $B$ the multiples of $3$. Then
\[
A=\{2,4,6,8,10\},\qquad B=\{3,6,9\},\qquad A\cap B=\{6\}.
\]
So $|A|=5$, $|B|=3$ and $|A\cap B|=1$. The union is $\{2,3,4,6,8,9,10\}$, which has $7$ elements. The sum $|A|+|B|=8$ is one too large, because $6$ was listed in both sets. Subtracting the overlap once repairs it:
\[
|A\cup B|=|A|+|B|-|A\cap B|=5+3-1=7.
\]
The same pattern, with one more set, is the pause below. The signs are the content of the theorem.

\begin{think}
How many integers in $\{1,\dots,100\}$ are divisible by $2$, $3$ or $5$? Adding $50+33+20$ counts some numbers twice or three times. Correct the count step by step, and watch the pattern of signs appear.
\end{think}

Here is the correction, before the general proof. There are $\lfloor 100/2\rfloor=50$ multiples of $2$, $\lfloor 100/3\rfloor=33$ multiples of $3$, and $\lfloor 100/5\rfloor=20$ multiples of $5$. Their sum is $103$. A number divisible by two of the three primes was counted twice. There are $\lfloor 100/6\rfloor=16$ multiples of $6$, $\lfloor 100/10\rfloor=10$ multiples of $10$, and $\lfloor 100/15\rfloor=6$ multiples of $15$. Subtract those:
\[
103-16-10-6=71.
\]
A multiple of $30$ was counted three times in the first sum, and then subtracted three times, once for each pair. It has disappeared from the running total. There are $\lfloor 100/30\rfloor=3$ such numbers, namely $30$, $60$ and $90$. Add them back:
\[
71+3=74.
\]
The signed sum is $50+33+20-16-10-6+3=74$. The signs are $+$, $+$, $+$, then $-$, $-$, $-$, then $+$. The theorem says that this is not an accident of the primes $2$, $3$ and $5$.

\begin{theorem}[inclusion--exclusion]\label{thm:incexc}
For finite sets $A_1,\dots,A_n$, writing $A_I=\bigcap_{i\in I}A_i$,
\[
\left|A_1\cup\dots\cup A_n\right|=\sum_{k=1}^n(-1)^{k+1}\sum_{|I|=k}|A_I|.
\]
\end{theorem}
\begin{proof}
The argument has three steps. First an identity of functions, one value at each element. Then the expansion of a product. Then a sum over the elements.

\emph{Setup.} For a set $A$, its \emph{indicator} $\mathbf 1_A$ is the function with $\mathbf 1_A(x)=1$ if $x\in A$ and $\mathbf 1_A(x)=0$ otherwise. Let $U=A_1\cup\dots\cup A_n$. Fix $x\in U$. Then $x$ lies in at least one $A_i$, so at least one factor $1-\mathbf 1_{A_i}(x)$ is zero, and
\[
1-\prod_{i=1}^n\bigl(1-\mathbf 1_{A_i}(x)\bigr)=1.
\]
The left side is also the indicator of $U$, evaluated at $x$. If $x$ lies in none of the sets, every factor equals $1$, the product equals $1$, and both sides are $0$. We only need the identity on $U$.

\emph{Expansion.} For numbers $z_1,\dots,z_n$,
\[
\prod_{i=1}^n(1+z_i)=\sum_{I\subseteq\{1,\dots,n\}}\prod_{i\in I}z_i,
\]
where the empty product, for $I=\emptyset$, equals $1$. Put $z_i=-\mathbf 1_{A_i}(x)$. Then
\[
\prod_{i=1}^n\bigl(1-\mathbf 1_{A_i}(x)\bigr)=\sum_{I}(-1)^{|I|}\mathbf 1_{A_I}(x).
\]
The term $I=\emptyset$ equals $1$. For a nonempty $I$, the product of the indicators is the indicator of the intersection $A_I$. Subtract the expansion from $1$:
\begin{align*}
1-\prod_{i=1}^n\bigl(1-\mathbf 1_{A_i}(x)\bigr)
&=1-\sum_{I}(-1)^{|I|}\mathbf 1_{A_I}(x)\\
&=1-\Bigl(1+\sum_{\emptyset\neq I}(-1)^{|I|}\mathbf 1_{A_I}(x)\Bigr)\\
&=\sum_{\emptyset\neq I}(-1)^{|I|+1}\mathbf 1_{A_I}(x).
\end{align*}
So, at every $x\in U$,
\[
\mathbf 1_U(x)=\sum_{\emptyset\neq I}(-1)^{|I|+1}\mathbf 1_{A_I}(x).
\]

\emph{Sum.} Sum both sides over $x\in U$. The left side adds one for each element of $U$, so it equals $|U|$. On the right, group the terms by the nonempty set $I$. The sum of $\mathbf 1_{A_I}$ over $x\in U$ equals $|A_I|$, because $A_I\subseteq U$. Grouping the sets $I$ by their size $k=|I|$ is the sum in the statement.
\end{proof}

For $n=2$ the theorem is the count we already checked: the $k=1$ terms are $|A_1|+|A_2|$, and the $k=2$ term is $-|A_1\cap A_2|$. For $n=1$ it says $|A_1|=|A_1|$. For the multiples of $2$, $3$ and $5$ up to $100$, it is the signed sum $74$.

\begin{figure}[ht]
\centering
\fig{ch08-venn}
\caption{Inclusion--exclusion for the multiples of $2$, $3$ and $5$ in $\{1,\dots,100\}$. The signed sum $50+33+20-16-10-6+3$ equals $74$, the size of the union. The innermost region holds three numbers, $30$, $60$ and $90$. The seven regions of the circles hold $27$, $14$, $7$, $13$, $7$, $3$ and $3$ numbers, and those seven counts add to $74$. Twenty-six integers lie in no circle: $1$, the $22$ primes from $7$ to $97$, and the composites $49$, $77$ and $91$.}
\label{fig:venn}
\end{figure}

The identity proved above is an identity of functions, $\mathbf 1_U=\sum_{\emptyset\neq I}(-1)^{|I|+1}\mathbf 1_{A_I}$, and the cardinalities are what you get by adding the values. In Chapter~\ref{ch:probability}, weighting the same identity by probabilities gives the probability of a union. Keeping only the first few blocks of terms also gives inequalities. That comparison is Exercise~\ref{ex:bonferroni} in Chapter~\ref{ch:probability}.

\section{Derangements}

A \emph{fixed point} of a permutation $\sigma$ is an element $i$ with $\sigma(i)=i$. A permutation with no fixed point is a \emph{derangement}. Write $D_n$ for the number of derangements of $\{1,\dots,n\}$. The opening list gave $D_3=2$. The empty set has one permutation, the empty function. There is no element that could be a fixed point, so $D_0=1$.

\begin{proposition}[derangements]\label{prop:derang}
For every integer $n\ge0$,
\[
D_n=n!\sum_{k=0}^n\frac{(-1)^k}{k!}.
\]
\end{proposition}
\begin{proof}
\emph{The sets.} For $n=0$ the sum is the single term $1$, and $0!\cdot 1=1=D_0$. Now take $n\ge 1$. Let $S$ be the set of all permutations of $\{1,\dots,n\}$, so $|S|=n!$ by Proposition~\ref{prop:counts}. Let $A_i$ be the set of permutations that fix $i$. A permutation lies in the union $A_1\cup\dots\cup A_n$ precisely when it has at least one fixed point. The derangements are the permutations in none of the sets $A_i$, so
\[
D_n=n!-\bigl|A_1\cup\dots\cup A_n\bigr|.
\]

\emph{The intersections.} Let $I\subseteq\{1,\dots,n\}$ with $|I|=k\ge 1$. A permutation in $A_I$ fixes every point of $I$, and the remaining $n-k$ points may be sent anywhere among themselves. There are $(n-k)!$ such permutations. There are $\binom nk$ sets $I$ of size $k$. Theorem~\ref{thm:incexc} therefore gives
\begin{align*}
\bigl|A_1\cup\dots\cup A_n\bigr|
&=\sum_{k=1}^n(-1)^{k+1}\binom nk(n-k)!\\
&=\sum_{k=1}^n(-1)^{k+1}\frac{n!}{k!},
\end{align*}
where the second line uses $\binom nk(n-k)!=n!/k!$.

\emph{Subtract.} Factor $n!$ out of the sum:
\[
D_n=n!-\,n!\sum_{k=1}^n\frac{(-1)^{k+1}}{k!}
=n!\Biggl(1-\sum_{k=1}^n\frac{(-1)^{k+1}}{k!}\Biggr).
\]
The expression in parentheses is the full sum from $k=0$. The $k=0$ term is $1$. For $k\ge 1$, the sign in front of $1/k!$ is $-(-1)^{k+1}$. And $-(-1)^{k+1}=(-1)^{k+2}=(-1)^k$. Thus
\[
D_n=n!\sum_{k=0}^n\frac{(-1)^k}{k!}.
\]
\end{proof}

Check $n=3$ against Figure~\ref{fig:hats}, using the count of the union rather than only the final number. The formula for the union is
\[
3!\Biggl(\frac{1}{1!}-\frac{1}{2!}+\frac{1}{3!}\Biggr)=6\Biggl(1-\frac12+\frac16\Biggr)=6\cdot\frac23=4.
\]
Four permutations of three hats have a fixed point, and the figure shows those four words. Then $D_3=6-4=2$, which matches $231$ and $312$. The same arithmetic for $n=4$ gives
\[
D_4=24\Biggl(1-1+\frac12-\frac16+\frac1{24}\Biggr)=24\cdot\frac{9}{24}=9.
\]
The first values are $D_0=1$, $D_1=0$, $D_2=1$, $D_3=2$ and $D_4=9$. For $n=1$ the only permutation is the identity, so $D_1=0$, and the sum is $1-1=0$. Exercise~\ref{ex:derangrec} finds a recurrence for $D_n$ without inclusion--exclusion, and it asks you to check the recurrence against these values.

The proportion of derangements is the finite sum $D_n/n!=\sum_{k=0}^n(-1)^k/k!$. Each larger $n$ adds one term. Whether the sums settle near a single number is a question about a sequence, and this chapter does not yet have a limit. Chapter~\ref{ch:sequences} says what a limit is. Chapter~\ref{ch:probability} returns to the proportion once that language exists.

\subsection*{Following a permutation around its cycles}
A permutation $\sigma$ of a finite set is a bijection from the set to itself. Starting at an element $a$, follow the list $a,\sigma(a),\sigma(\sigma(a)),\dots$. The set is finite, so some value repeats. The first repeated value must be $a$ itself. Suppose $\sigma^i(a)=\sigma^j(a)$ with $0\le i<j$, and suppose $i\ge 1$. Injectivity of $\sigma$ gives $\sigma^{i-1}(a)=\sigma^{j-1}(a)$, which is an earlier repetition. So the first repetition cannot have $i\ge 1$. It has $i=0$, and the repeated value is $a$. The distinct elements $a,\sigma(a),\dots$ up to that return form a \emph{cycle}.

Start again at an element that has not been visited, and you get another cycle. Two cycles that share an element share the successor of that element, and the predecessor, so they contain exactly the same elements. The cycles are therefore disjoint, and every element lies on the cycle that starts there. They partition the set. A fixed point is a cycle of length one.

The permutation with $1\mapsto 3\mapsto 2\mapsto 1$ and $4\mapsto 5\mapsto 4$ has cycles $(1\ 3\ 2)$ and $(4\ 5)$. It has five elements and two cycles, and no cycle of length one, so it is a derangement. Figure~\ref{fig:cycles} draws it. The picture makes the splice in Exercise~\ref{ex:derangrec} visible: a $2$-cycle can be lifted out, and a longer cycle can be cut. The same description prepares the cycle count of Exercise~\ref{ex:cycles}.

\begin{figure}[ht]
\centering
\fig{ch02-cycles}
\caption{One derangement of $\{1,2,3,4,5\}$, drawn as its cycles. The arrows follow $1\mapsto 3\mapsto 2\mapsto 1$ and $4\mapsto 5\mapsto 4$. There are two cycles and five elements. Neither cycle has length one, so no point is fixed.}
\label{fig:cycles}
\end{figure}

\section{Lattice paths}

The binomial coefficient has a second description. Let $a,b\in\N$. A \emph{lattice path} from $(0,0)$ to $(a,b)$ is a walk by unit steps east and north:
\[
\mathsf E\colon (x,y)\to(x+1,y),\qquad \mathsf N\colon (x,y)\to(x,y+1).
\]
Figure~\ref{fig:paths} draws one path to $(4,3)$, and a second picture that the Catalan exercise will use.

\begin{proposition}[counting paths]\label{prop:paths}
The number of lattice paths from $(0,0)$ to $(a,b)$ is $\binom{a+b}{a}$.
\end{proposition}
\begin{proof}
\emph{The word.} A path from $(0,0)$ to $(a,b)$ takes $a$ steps $\mathsf E$ and $b$ steps $\mathsf N$, in some order. It has $a+b$ steps in all. Reading the steps in order is a word of length $a+b$ on the alphabet $\{\mathsf E,\mathsf N\}$ with exactly $a$ letters $\mathsf E$. The path determines the word, and the word determines the path.

\emph{The positions.} The word is determined by the set of positions of the letter $\mathsf E$. That set is an $a$-element subset of $\{1,\dots,a+b\}$. Chapter~\ref{ch:sums} counts those subsets by $\binom{a+b}{a}$. Paths, words and $a$-element subsets are in bijection.

\emph{Check.} A path to $(4,3)$ has $4$ letters $\mathsf E$ and $3$ letters $\mathsf N$. The number is $\binom74=35$. Choosing the three positions of $\mathsf N$ instead gives $\binom73=35$, the same number. The word $\mathsf{ENEENEN}$ in Figure~\ref{fig:paths} is one of those $35$ paths.
\end{proof}

Several identities from Chapter~\ref{ch:sums} become one-line counts once the picture is drawn.

Pascal's rule is an arrival. A path to $(a,b)$, with $a\ge 1$ and $b\ge 1$, takes its last step either east, from $(a-1,b)$, or north, from $(a,b-1)$. The two sets of paths are disjoint. Proposition~\ref{prop:paths} turns the split into
\[
\binom{a+b}{a}=\binom{a+b-1}{a-1}+\binom{a+b-1}{a},
\]
which is Proposition~\ref{prop:binom}(a). If $a=0$, the only path to $(0,b)$ is a run of north steps, and there is one such path. Pascal's rule still holds numerically, because $\binom b0=\binom{b-1}{-1}+\binom{b-1}{0}$ and Chapter~\ref{ch:sums} sets $\binom{b-1}{-1}=0$. The case $b=0$ is the same with the roles exchanged.

Vandermonde's identity is a halfway point. Consider the paths from $(0,0)$ to $(r,\,m+n-r)$. There are $\binom{m+n}{r}$ of them. After $m$ steps the path is at some point $(k,\,m-k)$, where $k$ is the number of east steps so far. The first $m$ steps are a path to $(k,\,m-k)$, and there are $\binom mk$ of those. The remaining $n$ steps run from $(k,\,m-k)$ to $(r,\,m+n-r)$, so they contain $r-k$ east steps, and there are $\binom n{r-k}$ of them. Sum over $k$:
\[
\binom{m+n}{r}=\sum_{k}\binom mk\binom n{r-k}.
\]
Any term with $k>m$, or $k>r$, or $r-k>n$, is zero by the convention on binomial coefficients. For a check, take $m=n=r=2$. The sum is
\[
\binom20\binom22+\binom21\binom21+\binom22\binom20=1+4+1=6,
\]
and $\binom42=6$.

The hockey-stick identity sorts paths by their last east step. Count the paths to $(r+1,\,n-r)$. There are $\binom{n+1}{r+1}$ of them, since the total number of steps is $(r+1)+(n-r)=n+1$. Let the last east step run from $(r,h)$ to $(r+1,h)$, and suppose that every later step is north. Then $0\le h\le n-r$, because the path still has to climb from height $h$ to height $n-r$. Before that east step the walk is a path to $(r,h)$, and there are $\binom{r+h}{r}$ of those. After it, the climb is forced. Therefore
\[
\binom{n+1}{r+1}=\sum_{h=0}^{n-r}\binom{r+h}{r}.
\]
For $r=1$ and $n=3$ the right side is $\binom11+\binom21+\binom31=1+2+3=6$, and the left side is $\binom42=6$. The change of index $k=r+h$ rewrites the sum as $\sum_{k=r}^{n}\binom kr=\binom{n+1}{r+1}$. That is the hockey-stick identity of Exercise~\ref{ex:hockey}. Chapter~\ref{ch:sums} proves it by induction. The path count is a second proof.

\begin{figure}[ht]
\centering
\fig{ch01-paths}
\caption{Left: the word $\mathsf{ENEENEN}$ is one path from $(0,0)$ to $(4,3)$. It has four east steps and three north steps, so it is one of the $\binom74=\binom73=35$ paths. Right: a path from $(0,0)$ to $(4,4)$ that rises above the diagonal meets the line $y=x+1$. Reflecting the part after the first meeting point, drawn dotted and then in rust, produces a path to $(3,5)$. The good paths to $(4,4)$, those that stay on or below $y=x$, number $\binom84-\binom83=70-56=14$.}
\label{fig:paths}
\end{figure}

Exercise~\ref{ex:catalan} uses the right-hand picture to count the paths from $(0,0)$ to $(n,n)$ that never rise above the diagonal $y=x$. Those numbers are the \emph{Catalan numbers} $1,1,2,5,14,42,\dots$. The case drawn in the figure is $n=4$: there are $\binom84=70$ paths to $(4,4)$ in all, and $\binom83=56$ paths to $(3,5)$, so $70-56=14$ paths stay on or below the diagonal. Stanley lists more than two hundred other families counted by the same numbers. Chapter~\ref{ch:sums} already met them as Euler's count of triangulations of a polygon.

\section{Generating functions: keeping every count at once}
A recurrence explains the next term from the previous ones. A generating function stores the whole list. Nothing in this section converges or diverges. Each coefficient is a finite sum, and equality means that corresponding coefficients are equal.

\begin{think}
A strip of $n$ unit squares is tiled by tiles of lengths $1$ and $2$. What can the last tile be? Work out the counts for $n=0,1,2,3,4$ before writing a formula. Why must the empty strip have one tiling?
\end{think}

Let $u_n$ be the number of tilings. The empty strip has one tiling, the tiling that uses no tiles. That count has to be $1$. If it were $0$, the relation $u_2=u_1+u_0$ would give $u_2=1$, and the tiling made of one tile of length $2$ would be missing. So $u_0=1$. A strip of length $1$ has one tiling, a single tile of length $1$, and $u_1=1$. For length $2$ there are two tilings: two short tiles, or one long tile. For length $3$ the last tile is short, leaving a tiling of length $2$, or long, leaving a tiling of length $1$. That gives the three tilings short-short-short, short then long, and long then short. For length $4$ the same split gives $u_3+u_2=3+2=5$. The list begins $1,1,2,3,5$.

The two cases are the recurrence. For $n\ge 2$, remove the last tile. A last tile of length $1$ leaves a tiling of length $n-1$. A last tile of length $2$ leaves a tiling of length $n-2$. Every tiling arises once this way, and the two constructions can be undone by putting the tile back. Therefore $u_n=u_{n-1}+u_{n-2}$.

\begin{examplex}[a recurrence becomes an equation]
Let $U(t)=\sum_{n\ge 0}u_nt^n$. This is a \emph{formal power series}: $t$ is a marker for the length, and $[t^n]U(t)$ means the coefficient $u_n$. Multiplying by $t$ shifts the list one place to the right and inserts a zero at the start. Compare coefficients in $(1-t-t^2)U(t)$.

\emph{Degree $0$.} The constant term is $u_0=1$.

\emph{Degree $1$.} The coefficient is $u_1-u_0=1-1=0$.

\emph{Degree $n\ge 2$.} The coefficient is $u_n-u_{n-1}-u_{n-2}=0$, by the recurrence.

So $(1-t-t^2)U(t)=1$. We write $U(t)=1/(1-t-t^2)$ for this formal inverse: the equation $(1-t-t^2)U(t)=1$ determines every coefficient, and it determines them uniquely. The first coefficients are $1,1,2,3,5,8$. Exercise~\ref{ex:fibgcd} defines the Fibonacci sequence by $F_0=0$, $F_1=1$ and $F_{n+2}=F_{n+1}+F_n$. Then $u_n=F_{n+1}$, because $F_1=1$, $F_2=1$, and both sequences obey the same recurrence. The series organises the decomposition. The proof of the decomposition is still the case split on the last tile.
\end{examplex}

The same bookkeeping multiplies two series. If $A(t)=\sum_{n\ge 0}a_nt^n$ and $B(t)=\sum_{n\ge 0}b_nt^n$, the product is defined by collecting degree $n$:
\[
[t^n]\bigl(A(t)B(t)\bigr)=\sum_{j=0}^n a_jb_{n-j}.
\]
For each fixed $n$ the sum is finite. If $a_j$ counts objects of size $j$ and $b_\ell$ counts objects of a second kind of size $\ell$, then $a_jb_{n-j}$ counts ordered pairs of total size $n$ whose first object has size $j$. The sum over $j$ counts every such pair once. Multiplication is the product rule, together with a split by size. A product of any fixed finite number of series is the same idea, one factor at a time. A coefficient in a product of three series is a finite sum over triples, and grouping the triples in either order gives the same sum, so the product is associative.

One series will be used often enough to name it. Let $G(t)=1+t+t^2+\cdots=\sum_{n\ge 0}t^n$. Then
\[
(1-t)G(t)=1,
\]
because the constant term is $1$, and at every positive degree the coefficient is $1-1=0$. We write $G(t)=1/(1-t)$. For an integer $r\ge 1$, the power $1/(1-t)^r$ is the product of $r$ copies of $G$. Its coefficient of $t^n$ counts the $r$-tuples of nonnegative integers with sum $n$, by the multiplication rule just given. Proposition~\ref{prop:counts}(b) counts those tuples by stars and bars, with $n$ stars and $r-1$ bars, so
\[
[t^n]\frac1{(1-t)^r}=\binom{n+r-1}{r-1}\qquad(r\ge 1).
\]
For $r=2$ the right side is $\binom{n+1}{1}=n+1$. The pairs $(x,y)\in\N^2$ with $x+y=n$ are $(0,n),(1,n-1),\dots,(n,0)$, and there are $n+1$ of them. The fraction denotes a formal inverse. It does not ask for a numerical value of $t$.

For a finite set, an ordinary polynomial is enough. If each element of $\{1,\dots,m\}$ is either omitted or selected, and selecting $j$ contributes degree $j$, then
\[
\prod_{j=1}^m(1+t^j)
\]
has $[t^n]$ equal to the number of subsets with element sum $n$. The factor $1$ omits $j$. The factor $t^j$ selects it. For $m=3$ the polynomial is $(1+t)(1+t^2)(1+t^3)$. The subsets of $\{1,2,3\}$ with sum $3$ are $\{3\}$ and $\{1,2\}$, two subsets, and the coefficient of $t^3$ is $2$. This is the binomial theorem's expansion, with degree recording a sum of selected numbers rather than a count of selected factors.

\section{Workshop: remove the restriction by changing the objects}
How many $k$-element subsets of $\{1,\dots,n\}$ contain no consecutive integers? A formula is less useful than a transformation that explains why it holds.

\subsection*{Experiment and diagnosis}
Take $n=5$ and $k=2$. The admissible pairs are
\[
\{1,3\},\ \{1,4\},\ \{1,5\},\ \{2,4\},\ \{2,5\},\ \{3,5\}.
\]
There are six. The pair $\{1,2\}$ is missing because its entries differ by $1$, and so are $\{2,3\}$, $\{3,4\}$ and $\{4,5\}$. A product ``choose the first element, then choose the second'' does not have a fixed number of choices at the second stage: after $1$ there are three legal partners in this list, and after $3$ there is one. The product rule, in the form that needs a constant number of options, does not apply to that order of choosing. List a selected set in increasing order instead,
\[
1\le a_1<a_2<\cdots<a_k\le n,\qquad a_{i+1}\ge a_i+2.
\]
The extra space is compulsory. After every selected element except the last, at least one place stays unused. The aim is to remove exactly that compulsory space.

\subsection*{Mission and hints}
Find a bijection from admissible selections to ordinary $k$-subsets of a shorter interval. Treat the empty selection separately, and determine when there are no admissible selections at all.

\textbf{Hint 1.} The second selected element has one compulsory gap before it, the third has two. \textbf{Hint 2.} Subtract $i-1$ from the $i$th selected element. Check that this map is invertible before counting its target.

\subsection*{How to find it}
We want to replace a gap of at least two by a gap of at least one. If $b_i=a_i-(i-1)$, then
\[
b_{i+1}-b_i=(a_{i+1}-a_i)-1.
\]
The hypothesis $a_{i+1}-a_i\ge 2$ becomes $b_{i+1}-b_i\ge 1$. The simplest correction that removes one compulsory gap before the $i$th selected element is to subtract $i-1$. The upper endpoint decreases from $n$ to $n-k+1$. Figure~\ref{fig:gaps} shows one example: the selected positions $1,3,5$ in $\{1,\dots,7\}$ become $1,2,3$ in a shorter interval.

\begin{figure}[ht]
\centering
\fig{ch02-gaps}
\caption{Removing compulsory gaps. The three rust positions $1,3,5$ in $\{1,\dots,7\}$ have a forced unused place before each selection after the first. Subtracting $0,1,2$ sends them to the three charcoal positions $1,2,3$, an ordinary $3$-subset of $\{1,\dots,5\}$. The shorter interval has length $7-3+1=5$.}
\label{fig:gaps}
\end{figure}

\subsection*{Write-up}
For $k\ge 1$, an admissible set has $a_k\ge a_1+2(k-1)\ge 1+2k-2=2k-1$. Therefore none exists when $n<2k-1$. Now suppose $n\ge 2k-1$, and define $b_i=a_i-(i-1)$. Then $b_1=a_1\ge 1$, the gaps give $b_{i+1}\ge b_i+1$, and
\[
b_k=a_k-(k-1)\le n-k+1.
\]
So $1\le b_1<b_2<\cdots<b_k\le n-k+1$. Conversely, every strictly increasing list in this interval produces an admissible list by $a_i=b_i+(i-1)$. The two constructions undo each other: subtracting $i-1$ and then adding $i-1$ returns the original value. The desired count is therefore
\[
\binom{n-k+1}{k}\qquad(n\ge 2k-1).
\]
The upper index $n-k+1$ is at least $k$ in this range, so the binomial coefficient is the ordinary one from Chapter~\ref{ch:sums}. For the six pairs with $n=5$ and $k=2$, the formula gives $\binom{5-2+1}{2}=\binom42=6$.

For $k=0$ there is exactly one selection, the empty set. That value is needed when the counts for every $k$ are added.

Let $f_n$ be the number of all subsets of $\{1,\dots,n\}$ with no two consecutive elements, including the empty set. The case split in the first section gave $f_n=f_{n-1}+f_{n-2}$ for $n\ge 3$. Set $f_0=1$, counting the empty subset of the empty set. The same split at $n=2$ reads $f_2=f_1+f_0=3$, which matches the list, so the recurrence holds for every $n\ge 2$. Then $f_0=1=F_2$ and $f_1=2=F_3$, with the Fibonacci numbering of Exercise~\ref{ex:fibgcd}. Both sequences satisfy the same recurrence and the same two initial values, so $f_n=F_{n+2}$ for every $n\ge 0$. This is induction with two base cases, as in Chapter~\ref{ch:sums}.

Count the same subsets by size. The largest admissible $k$ is the largest integer with $2k-1\le n$, namely $k\le (n+1)/2$. Therefore
\[
F_{n+2}=\sum_{k=0}^{\lfloor(n+1)/2\rfloor}\binom{n-k+1}{k}.
\]
For $n=5$ the right side is
\[
\binom60+\binom51+\binom42+\binom33=1+5+6+1=13,
\]
and $F_7=13$. The bijection, the recurrence and the sum describe the same objects in three languages. Exercise~\ref{ex:count-circle} asks what changes on a circle, where the two endpoints can meet.

\begin{reminder}[choosing a counting method]
A local construction suggests a recurrence. A global restriction may suggest its complement. Several overlapping restrictions suggest inclusion--exclusion. If a constraint demands compulsory gaps, try removing those gaps by a bijection. In every method, check an empty object, the first nonempty object, and a case where no object is possible. A formula that misses one of these cases usually hides a wrong interpretation of the objects.
\end{reminder}

\section*{What this chapter prepares}
Counting, in this chapter, is a choice of language. The sum, product and bijection rules are the grammar. Inclusion--exclusion repairs an overlap, and the hat check is the overlap written as a formula: $D_n=n!\sum_{k=0}^n(-1)^k/k!$. Paths and gap-removal are bijections. A generating function is a list of finite counts, one coefficient at a time.

The proportion $D_n/n!$ is that same finite sum. As $n$ grows, the sum gains terms. Whether the numbers $D_n/n!$ approach one value is not a counting question. It is a question about a sequence. The next chapter, Sequences, Series and Limits, says what a limit is, and what it means to add infinitely many real numbers. The series $U(t)$ of this chapter stores counts as coefficients. The next chapter starts from sequences of real numbers, and only then adds those numbers up. The floor of a general real number is defined there as well.

The hat check will be read again as a probability, in Chapter~\ref{ch:probability}, after a limit has a meaning. This chapter stops at the exact count.

\section{Exercises}
\begin{exercise}[$\star$]\label{ex:stars}
Let $n\ge0$ and $k\ge1$ be integers. How many solutions in $\N^k$ does $x_1+\dots+x_k=n$ have? And in $(\N^*)^k$?
\end{exercise}

\begin{exercise}[$\star\star$]\label{ex:pigeon}
Let $n\ge1$. Choose $n+1$ distinct integers in $\{1,\dots,2n\}$. Show that two of them are coprime, and that one of them divides another.
\end{exercise}

\begin{exercise}[$\star\star$]\label{ex:phi}
Let $n\ge1$, and write $n=p_1^{a_1}\cdots p_r^{a_r}$ with distinct primes $p_i$. Using inclusion--exclusion, show that the number $\varphi(n)$ of integers in $\{1,\dots,n\}$ coprime to $n$ is
$n\prod_{i=1}^r\left(1-\frac1{p_i}\right)$.
\end{exercise}

\begin{exercise}[$\star\star$]\label{ex:derangrec}
\emph{Euler's recurrence.} Show directly (by a bijection, not from the formula) that the derangement numbers satisfy $D_n=(n-1)(D_{n-1}+D_{n-2})$ for $n\ge3$, and deduce $D_n=nD_{n-1}+(-1)^n$. Check both against the values in the text.
\end{exercise}

\begin{exercise}[$\star\star$]\label{ex:surj}
\emph{Surjections.} Let $k,n\ge0$. For these counts use $0^0=1$, because there is one map from the empty set to itself. Show that the number of surjections from a $k$-element set onto an $n$-element set is
\[
S(k,n)=\sum_{j=0}^n(-1)^j\binom nj(n-j)^k ,
\]
and deduce that $\sum_{j=0}^n(-1)^j\binom nj(n-j)^k=0$ for $k<n$ and $=n!$ for $k=n$. (Compare with Newton's forward differences, Theorem~\ref{thm:newton}.) How many surjections are there from $5$ elements onto $3$?
\end{exercise}

\begin{exercise}[$\star\star\star$]\label{ex:catalan}
\emph{The Catalan numbers.} Let $C_n$ be the number of paths from $(0,0)$ to $(n,n)$ with steps $\mathsf E,\mathsf N$ that never rise above the diagonal $y=x$.
\begin{enumerate}[label=(\alph*)]
\item For $n\ge1$, a path to $(n,n)$ that does rise above the diagonal has a first point on the line $y=x+1$. Reflect the rest of the path in that line, as in Figure~\ref{fig:paths}. Show that this gives a bijection between the ``bad'' paths to $(n,n)$ and all paths to $(n-1,n+1)$.
\item Deduce $C_n=\binom{2n}n-\binom{2n}{n-1}=\frac1{n+1}\binom{2n}n$, and compute $C_1,\dots,C_6$.
\item Show that $C_n$ also counts the ways to write $n$ pairs of parentheses correctly ($n=3$: ((())), (()()), (())(), ()(()), ()()()), and that $C_{n+1}=\sum_{k=0}^nC_kC_{n-k}$.
\end{enumerate}
\end{exercise}

\begin{exercise}[$\star\star\star$]\label{ex:erdos}
\emph{Erd\H{o}s--Szekeres.} Show that every sequence of $n^2+1$ distinct real numbers has a monotone subsequence of length $n+1$.
\end{exercise}
\begin{figure}[ht]
\centering
\fig{ch08-erdos}
\caption{Ten points $(i,a_i)$ for the sequence $3,2,1,6,5,4,9,8,7,10$. Under each point is the pair $(u_i,d_i)$ from Exercise~\ref{ex:erdos}: the lengths of the longest increasing and longest decreasing subsequences ending at that term. The ten pairs are all different, so they do not fit in the nine boxes of $\{1,2,3\}^2$. Here $u_{10}=4$, and the increasing subsequence $3,6,9,10$ is drawn in ochre. The shorter sequence $3,2,1,6,5,4,9,8,7$ has nine terms and no monotone subsequence of length $4$, so $n^2$ terms are not enough when $n=3$.}
\label{fig:erdos}
\end{figure}

\begin{exercise}[$\star$]\label{ex:count-coefficient}
Find the coefficient of $t^6$ in $(1+t^2+t^4)(1+t+t^2+t^3)(1+t^3)$. Interpret it as a count of triples $(a,b,c)$ with $a\in\{0,2,4\}$, $b\in\{0,1,2,3\}$, $c\in\{0,3\}$ and $a+b+c=6$. Explain why the product rule alone does not give the answer.
\end{exercise}

\begin{exercise}[$\star\star$]\label{ex:count-tiles}
Let $v_n$ count tilings of a strip of length $n$ using tiles of lengths $1$ and $3$. Prove a recurrence with its initial conditions, obtain its formal generating function, and compute $v_6$ in two ways: by recurrence and by classifying tilings according to their number of length-$3$ tiles.
\end{exercise}

\begin{exercise}[$\star\star$]\label{ex:count-derang-gf}
For the derangement numbers define the formal series $E(t)=\sum_{n\ge0}(D_n/n!)t^n$. Using Proposition~\ref{prop:derang}, show coefficientwise that
\[
E(t)=\left(\sum_{j\ge0}\frac{(-1)^jt^j}{j!}\right)\frac1{1-t}.
\]
Do not substitute a value of $t$. Deduce $D_n=nD_{n-1}+(-1)^n$ for $n\ge1$ directly from this product identity. Compare with the bijective proof in Exercise~\ref{ex:derangrec}.
\end{exercise}

\begin{exercise}[$\star\star\star$]\label{ex:count-circle}
Place the labels $1,\dots,n$ around a circle, with $n\ge3$. Two labels are adjacent if consecutive around the circle, so $1$ and $n$ are adjacent too. For $1\le k\le\lfloor n/2\rfloor$, prove that the number of $k$-subsets containing no adjacent labels is
\[
\binom{n-k}{k}+\binom{n-k-1}{k-1}
=\frac{n}{n-k}\binom{n-k}{k}.
\]
Treat $k=0$ and $k>\lfloor n/2\rfloor$ separately. The labels are fixed: rotations of a chosen subset are counted as different subsets unless they give exactly the same set.
\end{exercise}

\begin{beauty}
One restriction can be read in more than one language. A compulsory gap becomes a bijection. A last tile becomes a recurrence. A recurrence becomes a generating function, which here is a list of counts. Inclusion--exclusion turns the integers coprime to $n$ into the product of Exercise~\ref{ex:phi}. The surjection count of Exercise~\ref{ex:surj} is a finite difference, the same kind of difference Chapter~\ref{ch:sums} used for polynomials.

The numbers $D_n/n!$ form a sequence. The next chapter says what it means for a sequence to have a limit. Chapter~\ref{ch:probability} will then place a probability on a finite set of permutations. The proportion of derangements is waiting for that limit.
\end{beauty}

\section{Hints}
\begin{hint}{ex:stars}
Stars and bars. For positive solutions, set $y_i=x_i-1$.
\end{hint}

\begin{hint}{ex:pigeon}
Boxes for coprimality: $\{1,2\},\{3,4\},\dots$. For divisibility: write each integer as $2^a m$ with $m$ odd, and put it in the box of $m$.
\end{hint}

\begin{hint}{ex:phi}
Let $A_i$ be the multiples of $p_i$ in $\{1,\dots,n\}$. Then $|A_I|=n/\prod_{i\in I}p_i$, and expand $\prod\left(1-\frac1{p_i}\right)$.
\end{hint}

\begin{hint}{ex:derangrec}
In a derangement, $n$ is sent to some $k\ne n$ ($n-1$ choices). Either $k$ is sent back to $n$, or not; in the second case, pretend that $k$'s image is ``what $n$ would have had''.
\end{hint}

\begin{hint}{ex:surj}
Let $A_j$ be the set of maps missing the $j$th element of the target; count the maps in no $A_j$.
\end{hint}

\begin{hint}{ex:catalan}
(a) The reflection $(x,y)\mapsto(y-1,x+1)$ swaps $\mathsf E$ and $\mathsf N$ steps and sends $(n,n)$ to $(n-1,n+1)$; every path to $(n-1,n+1)$ must cross the line $y=x+1$. (c) A good path is $\mathsf E$, a good path, $\mathsf N$, a good path: split at the first return to the diagonal.
\end{hint}

\begin{hint}{ex:erdos}
Attach to each term the lengths of the longest increasing and decreasing subsequences ending there. Show that two different terms get different pairs.
\end{hint}

\begin{hint}{ex:count-coefficient}
Split according to $c=0$ or $c=3$, then test the three allowed values of $a$. The total-degree condition couples the choices.
\end{hint}
\begin{hint}{ex:count-tiles}
Remove the last tile. For the second count, a tiling with $j$ length-$3$ tiles contains $6-3j$ length-$1$ tiles, hence $6-2j$ tiles in all.
\end{hint}
\begin{hint}{ex:count-derang-gf}
Multiply by $1-t$ and compare the coefficient of $t^n$. The coefficient coming from $tE(t)$ is $D_{n-1}/(n-1)!$.
\end{hint}
\begin{hint}{ex:count-circle}
Split according to whether label $n$ is selected. In the second case, remove $n$ and its two neighbours, then apply the workshop's line formula to the remaining interval.
\end{hint}

\section{Solutions}
\begin{solution}{ex:stars}
\idea A solution of $x_1+\dots+x_k=n$ in $\N^k$ is a multiset of size $n$ on $k$ labels ($x_i$ = how many times label $i$ is used): this is stars and bars. For positive solutions, shift each variable by $1$ to fall back on the first case.
\writeup A solution is a row of $n$ stars split by $k-1$ bars into $k$ groups (possibly empty), so by Proposition~\ref{prop:counts}(b) there are $\binom{n+k-1}{n}=\binom{n+k-1}{k-1}$. For positive solutions, put $y_i=x_i-1$:
$y_1+\dots+y_k=n-k$ in $\N^k$. There are no solutions if $n<k$; for $n\ge k$ the count is $\binom{n-1}{k-1}$.
\end{solution}

\begin{solution}{ex:pigeon}
\idea ``Among $n+1$ objects, two share a property'' is pigeonhole with $n$ boxes; the art is to design the boxes. For coprimality, the boxes should contain coprime pairs: consecutive integers are coprime. For divisibility, the boxes should be chains under divisibility: $\{m,2m,4m,\dots\}$ with $m$ odd, and there are exactly $n$ odd numbers available.
\writeup Split $\{1,\dots,2n\}$ into the $n$ pairs $\{2k-1,2k\}$. Two of the $n+1$ chosen numbers lie in the same pair; they are consecutive, hence coprime (a common divisor divides their difference $1$).
For divisibility, write each chosen number as $2^am$ with $m$ odd. There are only $n$ odd numbers in $\{1,\dots,2n\}$, so two chosen numbers $2^am$ and $2^bm$ share the same $m$;
the one with the smaller power of $2$ divides the other.
\end{solution}

\begin{solution}{ex:phi}
\idea Counting elements that avoid several properties is the complement of a union, which inclusion--exclusion computes from the sizes of intersections. Here the intersections are easy: being divisible by several distinct primes means being divisible by their product. Finally, recognise the alternating sum as an expanded product.
\writeup If $n=1$, the empty product gives $1$, as it should: $\gcd(1,1)=1$. For $n>1$, let $A_i$ be the set of multiples of $p_i$ in $\{1,\dots,n\}$. For $I\subset\{1,\dots,r\}$, $A_I$ is the set of multiples of $\prod_{i\in I}p_i$ (distinct primes all divide $x$ iff their product does), so $|A_I|=n/\prod_{i\in I}p_i$.
An integer is coprime to $n$ if and only if it lies in no $A_i$, so by inclusion--exclusion
\[
\begin{aligned}\varphi(n)&=n-\sum_{\emptyset\ne I}(-1)^{|I|+1}\frac n{\prod_{i\in I}p_i}\eqbr=n\sum_{I}\prod_{i\in I}\left(-\frac1{p_i}\right)\eqbr=n\prod_{i=1}^r\left(1-\frac1{p_i}\right).\end{aligned}
\]
\end{solution}

\begin{solution}{ex:derangrec}
\idea Classify derangements by what happens to $n$: it goes to some $k$, and the question is whether $k$ goes back to $n$ (then the pair is a $2$-cycle and the rest is a smaller derangement) or not (then $n$ can be ``spliced out'' of its cycle, leaving a derangement of $n-1$ elements, and the splice can be undone if we remember $k$). The second recurrence is a one-line consequence once one looks at $D_n-nD_{n-1}$.
\writeup Let $\sigma$ be a derangement of $\{1,\dots,n\}$ and $k=\sigma(n)\ne n$: $n-1$ choices. \emph{Case 1}: $\sigma(k)=n$. Then $\sigma$ swaps $k$ and $n$, and restricted to the other $n-2$ elements it is a derangement; conversely any derangement of those elements extends to one of this kind. Count: $D_{n-2}$ for each $k$. \emph{Case 2}: $\sigma(k)\ne n$. Let $i=\sigma^{-1}(n)$, so $i\ne k$ and $i\ne n$. Define $\tau$ on $\{1,\dots,n-1\}$ by $\tau(i)=k$ and $\tau(j)=\sigma(j)$ for $j\ne i$ (these $\sigma(j)$ are all $\ne n$ since $j\ne i$, and $\ne k$ since $k=\sigma(n)$ and $j\ne n$). Then $\tau$ is a permutation of $\{1,\dots,n-1\}$ without fixed point: $\tau(j)=\sigma(j)\ne j$ for $j\ne i$, and $\tau(i)=k\ne i$. Conversely, given $k$ and a derangement $\tau$ of $\{1,\dots,n-1\}$, set $i=\tau^{-1}(k)$, $\sigma(n)=k$, $\sigma(i)=n$ and $\sigma=\tau$ elsewhere: a derangement with $\sigma(k)=\tau(k)\ne n$. The two constructions are inverse, so case~2 contributes $D_{n-1}$ for each $k$. Hence $D_n=(n-1)(D_{n-1}+D_{n-2})$.

For the second form, put $E_n=D_n-nD_{n-1}$. Then $E_n=(n-1)D_{n-1}+(n-1)D_{n-2}-nD_{n-1}=-(D_{n-1}-(n-1)D_{n-2})=-E_{n-1}$, and $E_2=D_2-2D_1=1$, so $E_n=(-1)^n$. Check: $D_4=3(2+1)=9=4\cdot2+1$, $D_5=4(9+2)=44=5\cdot9-1$.
\end{solution}

\begin{solution}{ex:surj}
\idea ``Surjective'' means ``misses no element'', the complement of a union of $n$ events ``misses element $j$'', and the intersections of those events are easy to count: a map missing a given set of $j$ elements is a map into the other $n-j$. Inclusion--exclusion does the rest. The vanishing for $k<n$ is then a combinatorial fact (there are no surjections) proved by an algebraic identity, or the reverse: both readings are instructive.
\writeup Let $A_j$ be the set of maps $\{1,\dots,k\}\to\{1,\dots,n\}$ whose image misses $j$. For $J$ of size $j$, $A_J$ is the set of maps into the $n-j$ elements outside $J$: $|A_J|=(n-j)^k$. The surjections are the maps in no $A_j$, so by Theorem~\ref{thm:incexc}
\[
S(k,n)=n^k-\sum_{j=1}^n(-1)^{j+1}\binom nj(n-j)^k=\sum_{j=0}^n(-1)^j\binom nj(n-j)^k .
\]
For $k<n$ there is no surjection (the image has at most $k<n$ elements), so the sum is $0$; for $k=n$ a surjection between sets of the same size is a bijection, so the sum is $n!$. In the language of Chapter~\ref{ch:sums}, the sum is the $n$th forward difference of the polynomial $x^k$ evaluated at $0$, which vanishes for $\deg<n$ and equals $n!$ times the leading coefficient for $\deg=n$: Theorem~\ref{thm:newton} and Proposition~\ref{prop:counts} say the same thing. From $5$ elements onto $3$: $3^5-3\cdot2^5+3\cdot1^5=243-96+3=150$.
\end{solution}

\begin{solution}{ex:catalan}
\idea The reflection trick converts a hard count (paths avoiding a region) into two easy counts (all paths, minus the bad ones), by matching the bad paths bijectively with paths to a shifted endpoint. The key observations are: a bad path has a well-defined \emph{first} bad moment; reflection in the line $y=x+1$ maps lattice paths to lattice paths (it swaps the two kinds of steps) and moves the endpoint; and every path to the shifted endpoint is bad, so the map can be inverted by the same recipe. For (c), think of $\mathsf E$ as ``('' and $\mathsf N$ as ``)''; the recurrence comes from splitting a path at its first return to the diagonal, where the first step is $\mathsf E$ and the matching step is $\mathsf N$.
\writeup The empty path gives $C_0=1$, consistent with the formula in (b). For the reflection argument take $n\ge1$. (a) Let $\gamma$ be a path to $(n,n)$ that rises above the diagonal: at some moment its height exceeds its abscissa, and since the coordinates change by $1$ at a time, there is a first point $Q$ of $\gamma$ on the line $y=x+1$. The reflection $\rho:(x,y)\mapsto(y-1,x+1)$ in this line maps an $\mathsf E$ step to an $\mathsf N$ step and conversely. Replace the part of $\gamma$ after $Q$ by its image under $\rho$: the result $\gamma'$ is a lattice path from $(0,0)$ to $\rho(n,n)=(n-1,n+1)$, whose first point on the line $y=x+1$ is still $Q$ (the part before $Q$ is unchanged). Conversely, any path $\delta$ to $(n-1,n+1)$ ends above the line, so it has a first point $Q$ on it; reflecting the part after $Q$ gives a path to $(n,n)$ which touches the line at $Q$, hence is bad, and whose first point on the line is $Q$. The two constructions undo each other, so they are inverse bijections.

(b) By (a) and Proposition~\ref{prop:paths}, the bad paths number $\binom{2n}{n-1}$, so $C_n=\binom{2n}n-\binom{2n}{n-1}$. Now
$\binom{2n}{n-1}=\binom{2n}n\cdot\frac{n}{n+1}$ (write both with factorials), so $C_n=\binom{2n}n\left(1-\frac n{n+1}\right)=\frac1{n+1}\binom{2n}n$.
Values: $C_1=1$, $C_2=2$, $C_3=5$, $C_4=14$, $C_5=42$, $C_6=132$.

(c) Read $\mathsf E$ as an opening parenthesis and $\mathsf N$ as a closing one. A word of $n$ pairs is correctly parenthesised exactly when every prefix has at least as many opening as closing parentheses, i.e.\ the path never rises above the diagonal. For the recurrence, a non-empty good path to $(n+1,n+1)$ starts with $\mathsf E$ and returns to the diagonal for the first time at some point $(k+1,k+1)$, $0\le k\le n$, by a step $\mathsf N$. Between these two steps lies a good path translated to start at $(1,0)$ and end at $(k+1,k)$, which stays strictly below the diagonal, i.e.\ a good path of length $k$ in the shifted frame; after the return lies a good path of length $n-k$. Conversely any such pair gives a good path. So $C_{n+1}=\sum_{k=0}^nC_kC_{n-k}$, with $C_0=1$: $C_4=C_0C_3+C_1C_2+C_2C_1+C_3C_0=5+2+2+5=14$.
\end{solution}

\begin{solution}{ex:erdos}
\idea $n^2+1$ and $n^2$ suggest pigeonhole with $n^2$ boxes, the pairs in $\{1,\dots,n\}^2$. So label each term with a pair of numbers that are at most $n$ if there is no long monotone subsequence, and show the labels are distinct. The hint's pair (longest increasing, longest decreasing ending here) works because a later term can always extend one of the two subsequences of an earlier term.
\writeup For each term $a_i$, let $u_i$ (resp.\ $d_i$) be the length of the longest increasing (resp.\ decreasing) subsequence ending at $a_i$. If $i<j$ then $(u_i,d_i)\ne(u_j,d_j)$:
if $a_i<a_j$, appending $a_j$ to an increasing subsequence ending at $a_i$ gives $u_j\ge u_i+1$; if $a_i>a_j$, then $d_j\ge d_i+1$ likewise. If all the $u_i,d_i$ were at most $n$, the $n^2+1$ distinct pairs would lie in
$\{1,\dots,n\}^2$, which has $n^2$ elements: impossible by pigeonhole. So some $u_i$ or $d_i$ is at least $n+1$.
Chapter~\ref{ch:orders} proves the same bound as a case of Dilworth's theorem, and the sequence of $n$ decreasing blocks of length $n$ shows that $n^2$ terms are not enough.
\end{solution}

\begin{solution}{ex:count-coefficient}
\idea A coefficient of a product records all ways to split its degree among the factors. The small allowed sets make listing by the last component shorter than expanding the whole polynomial.
\writeup For $c=0$ we need $a+b=6$, and only $(a,b)=(4,2)$ is allowed. For $c=3$ we need $a+b=3$, giving $(a,b)=(0,3)$ and $(2,1)$. Thus the coefficient is $3$. The unrestricted product rule gives $3\cdot4\cdot2=24$ triples of all sums; it cannot impose the shared equation $a+b+c=6$ without a further decomposition.
\end{solution}

\begin{solution}{ex:count-tiles}
\idea Removing the last tile partitions all tilings into two disjoint, reversible constructions. This yields the recurrence. Fixing the number of long tiles instead gives a word with a specified number of each type, hence a binomial coefficient.
\writeup There is one tiling each of lengths $0$, $1$ and $2$, so $v_0=v_1=v_2=1$. For $n\ge3$, the last tile has length $1$ or $3$, leaving a tiling of length $n-1$ or $n-3$ respectively. Hence $v_n=v_{n-1}+v_{n-3}$. If $V(t)=\sum_{n\ge0}v_nt^n$, coefficient comparison gives $(1-t-t^3)V(t)=1$, or $V(t)=1/(1-t-t^3)$ as a formal identity. The recurrence gives $v_3=2$, $v_4=3$, $v_5=4$, $v_6=6$. Alternatively, for $j=0,1,2$ long tiles there are $6-3j$ short tiles, and choosing which of the $6-2j$ tile positions hold long tiles gives $\binom{6-2j}{j}$. The total is $\binom60+\binom41+\binom22=1+4+1=6$. The empty tiling is necessary for the recurrence at $n=3$.
\end{solution}

\begin{solution}{ex:count-derang-gf}
\idea The coefficient in the finite derangement formula is a partial sum. Multiplication by $1+t+t^2+\cdots$ turns coefficients into partial sums, and multiplication back by $1-t$ takes their consecutive differences.
\writeup The coefficient of $t^n$ on the proposed right side is $\sum_{j=0}^n(-1)^j/j!=D_n/n!$, proving equality of every coefficient. Multiplication by $1-t$ gives
\[
(1-t)E(t)=\sum_{j\ge0}\frac{(-1)^jt^j}{j!}.
\]
For $n\ge1$ its $t^n$ coefficient is $D_n/n!-D_{n-1}/(n-1)!=(-1)^n/n!$. Multiplying by $n!$ gives the claimed recurrence. At degree zero the identity says $D_0=1$. Only finite coefficient sums were used; no convergence or interchange of numerical infinite sums is involved. The bijection of Exercise~\ref{ex:derangrec} explains the underlying objects; this calculation explains the alternating correction term.
\end{solution}

\begin{solution}{ex:count-circle}
\idea A circle has no preferred endpoint, but fixing one label creates a two-case partition. If the label is selected, its neighbours are forbidden; if it is omitted, the remaining labels form a line. The workshop already counted both kinds of line selections.
\writeup Suppose $1\le k\le\lfloor n/2\rfloor$. If $n$ is omitted, choose $k$ nonconsecutive elements from the line $1,\dots,n-1$, giving $\binom{n-k}{k}$. If $n$ is selected, neither $1$ nor $n-1$ is selected. We choose the other $k-1$ labels from $2,\dots,n-2$, a line of length $n-3$. For $k=1$ this is the unique empty selection; for $k\ge2$ the workshop gives $\binom{(n-3)-(k-1)+1}{k-1}=\binom{n-k-1}{k-1}$. The same expression equals $1$ at $k=1$. These cases are disjoint and exhaustive. Since
\[
\binom{n-k-1}{k-1}=\frac{k}{n-k}\binom{n-k}{k},
\]
the sum has the required form. For $k=0$ there is one empty subset. If $k>\lfloor n/2\rfloor$, map each selected label to its clockwise neighbour. These neighbours are distinct and unselected, so $k\le n-k$ would be necessary; no such subset exists. This final argument also checks the feasibility boundary independently of the formula.
\end{solution}

\section*{Sources and further reading}
\begin{itemize}\raggedright
\item R.~L.~Graham, D.~E.~Knuth and O.~Patashnik, \emph{Concrete Mathematics}, 2nd ed., Addison-Wesley, 1994, Chapters~5 and~7. Binomial coefficients and generating functions.
\item R.~P.~Stanley, \emph{Catalan Numbers}, Cambridge University Press, 2015. Many families counted by the Catalan numbers, with their history.
\item M.~Renault, \emph{Lost (and found) in translation: Andr\'e's actual method and its application to the generalized ballot problem}, Amer.\ Math.\ Monthly 115 (2008), 358--363. The distinction between Andr\'e's own bijection and the reflection proof used here.
\item P.~R.~de Montmort, \emph{Essay d'analyse sur les jeux de hazard}, Paris, 1708; 2nd ed.\ 1713, with the correspondence with Nicolaus Bernoulli on the \emph{probl\`eme des rencontres}. L.~Euler, \emph{Calcul de la probabilit\'e dans le jeu de rencontre}, M\'emoires de l'Acad\'emie des sciences de Berlin 7 (1751, printed 1753), 255--270.
\item P.~Erd\H{o}s and G.~Szekeres, \emph{A combinatorial problem in geometry}, Compositio Math.\ 2 (1935), 463--470. The monotone-subsequence theorem.
\item P.~G.~L.~Dirichlet, \emph{Verallgemeinerung eines Satzes aus der Lehre von den Kettenbr\"uchen}, Bericht Berlin Akad.\ (1842), 93--95. Pigeonhole applied to rational approximation.
\end{itemize}
